% ==========================================
% 7. Conclusion
% ==========================================
%
% Status: written last, together. About one page.
%
% Principle: the conclusion summarizes; it introduces no new result, no new
% argument and (as a rule) no new citation. It answers "what did we learn, and
% what remains to be done".

\section{Conclusion}
\label{sec:conclusion}

% ------------------------------------------
% Summary of the study
% ------------------------------------------
% TODO [LAST]: We recall in two or three sentences the problem and the approach,
%              without repeating the introduction.

% ------------------------------------------
% Main findings
% ------------------------------------------
% TODO [LAST]: We state the answer to each research question in one or two
%              sentences, with the key quantified result, in the same order as
%              in Section~\ref{sec:discussion-rq}.

% ------------------------------------------
% Contributions
% ------------------------------------------
% TODO [LAST]: We restate the contributions announced in the introduction
%              (Section~\ref{sec:introduction}) and verify that each one is
%              backed by a result.

% ------------------------------------------
% Limitations
% ------------------------------------------
% TODO [LAST]: We summarize in a few sentences the main limitations from
%              Section~\ref{sec:discussion-limits}.

% ------------------------------------------
% Future work
% ------------------------------------------
% TODO [LAST]: We propose concrete follow-up work (experiments, improvements,
%              open questions), each linked to a limitation or to an unexpected
%              result.

% ------------------------------------------
% Final checks
% ------------------------------------------
% TODO [CHECK]: The conclusion answers exactly the research questions asked in
%               Section~\ref{sec:rq-hypotheses}, no more, no less.
% TODO [CHECK]: The abstract, the introduction and the conclusion are reread
%               side by side and tell the same story.
% TODO [CHECK]: The final report is read end to end by someone who did not write
%               it, and the remarks are handled before submission.
